%=====================================================================
% resume-plain.tex —— Plain TeX 简历示例
% 编译方式（任选其一）：
%   pdftex resume-plain.tex   → 直接产出 resume-plain.pdf（推荐）
%   tex     resume-plain.tex  → 产出 resume-plain.dvi，再 dvipdfmx 转换
% 说明：纯 Knuth plain 格式，不依赖 LaTeX / \documentclass / 任何宏包。
%       字体使用 Computer Modern（TeX 发行版自带）。
%=====================================================================

%------------------------- 全局设置 -------------------------
% 中文排版说明：本文件按引擎分支设置字体。
%   XeTeX     → 系统字体（Songti SC），中文可显示；
%   pdfTeX    → 默认 Computer Modern，中文不可见（本文档的历史行为）；
%   NTex      → 中文字体 FandolSong（由前端 set_otf_font 注入，见下），
%               并开启 \utfinputmode=1 以便源文件直写中文。
% 引擎判别用 \ifx\utfinputmode\undefined——\utfinputmode 是 NTex 独有的
% 内部整数参数（M9 中文刀 2），pdfTeX/XeTeX 上未定义，故能精确分流。
\hsize=17cm                     % 版心宽度
\vsize=24.5cm                   % 版心高度（压缩页边距，一页放下）
\parindent=0pt                  % 不缩进
\baselineskip=14pt
\ifx\utfinputmode\undefined
  % ---------- 其它引擎（pdfTeX / XeTeX / LuaTeX） ----------
  \ifx\XeTeXrevision\undefined
    \font\sectfont=cmr12 scaled \magstep1     % pdfTeX：Computer Modern
  \else
    \XeTeXlinebreaklocale "zh"                 % 中文断行规则
    \XeTeXlinebreakskip = 0pt plus 1pt minus .1pt
    \font\tenrm="Songti SC" at 11pt            % 正文（宋体）
    \font\tenbf="Songti SC/B" at 11pt          % 粗体
    \font\tentt="Menlo" at 10.5pt              % 等宽
    \font\sectfont="Songti SC/B" at 14pt       % 节标题
    \tenrm                                     % 关键：激活中文字体为当前字体
    \baselineskip=16pt
  \fi
\else
  % ---------- NTex（含 Tauri / WASM 实时预览） ----------
  \utfinputmode=1                              % UTF-8 直写：多字节 → 单个字符 token
  % Fandol Song（CTAN，GPL）由前端 fetch 后经 set_otf_font 注入；
  % 文件主名即 TeX 字体名。中文暂无粗体件，\bf 与正文同轮廓（已知简化）。
  % 每个字体单独试加载：失败则 \ifx<cs>\nullfont 成立（TeX 的字体标识符比较），
  % 退回 Computer Modern——中文仍缺字，但西文至少可读（可恢复错误，不中断作业）。
  \font\tenrm=FandolSong-Regular at 11pt
  \ifx\tenrm\nullfont \font\tenrm=cmr10 \fi
  \font\tenbf=FandolSong-Regular at 11pt
  \ifx\tenbf\nullfont \font\tenbf=cmbx10 \fi
  \font\sectfont=FandolSong-Regular at 14pt
  \ifx\sectfont\nullfont \font\sectfont=cmr12 scaled \magstep1 \fi
  \tenrm
  \baselineskip=16pt
  % NTex 的内嵌 plain 尚未装配数学字体族（属 C 档 `.fmt` 路线；见
  % docs/KNOWN-SIMPLIFICATIONS.md）。不装配时 fam≠0 的 \mathchar 会落到
  % 正文字体上——本文档的 \bulchar（CMSY 0F，见下方 \mathchardef）就会报
  % `Missing character: no ^^F in font FandolSong-Regular`。此处显式装配
  % 三族（文字/数学斜体/符号/大算符），全部走内嵌 CM TFM。
  \font\tenmrm=cmr10 \font\sevenmrm=cmr7 \font\fivemrm=cmr5
  \font\tenmi=cmmi10 \font\sevenmi=cmmi7 \font\fivemi=cmmi5
  \font\tensy=cmsy10 \font\sevensy=cmsy7 \font\fivesy=cmsy5
  \font\tenex=cmex10
  \textfont0=\tenmrm \scriptfont0=\sevenmrm \scriptscriptfont0=\fivemrm
  \textfont1=\tenmi  \scriptfont1=\sevenmi  \scriptscriptfont1=\fivemi
  \textfont2=\tensy  \scriptfont2=\sevensy  \scriptscriptfont2=\fivesy
  \textfont3=\tenex  \scriptfont3=\tenex   \scriptscriptfont3=\tenex
\fi
% 等宽字体（\tt）保持 Computer Modern Typewriter：文件里只用于 URL 等
% 纯 ASCII 内容，中文不出现在等宽语境。

%------------------------- 自定义宏 -------------------------
% 节标题：粗体大字 + 横线
\newdimen\secgap  \secgap=6pt
\def\section#1{%
  \vskip\secgap
  \hrule height .8pt
  \vskip 4pt
  {\sectfont #1}%
  \vskip 4pt
}

% 单行条目：左粗体标题，右日期（同一行两端对齐）
% 用法：\entryline{公司 / 职位}{2024.07 -- 至今}
\def\entryline#1#2{\line{\bf #1\hfil #2}\smallskip}

% 描述行 / 列表项
% 描述行 / 列表项（bullet = CMSY 字体 0F 槽位，\mathchar"220F）
\mathchardef\bulchar="220F
\def\item#1{\line{\hbox to 1.2em{\hss$\bulchar$\hss}\vtop{\hsize=15.4cm\noindent #1}}\smallskip}

% 联系方式一行内的分隔
\def\sep{\quad$\cdot$\quad}

%------------------------- 正文开始 -------------------------

\centerline{\bf 张\enspace 三}
\medskip
\centerline{北京 \sep 138-0000-0000 \sep zhangsan@example.com}
\centerline{\tt github.com/zhangsan \sep blog.zhangsan.dev}
\medskip

%========================== 教育背景 ==========================
\section{教育背景}

\entryline{某某大学\quad 计算机科学与技术\quad 硕士}{2022.09 -- 2025.06}
\item{GPA 3.8/4.0；研究方向：编译器与程序分析。}
\entryline{某某大学\quad 软件工程\quad 本科}{2018.09 -- 2022.06}
\item{GPA 3.6/4.0，专业排名前 5\%；获国家奖学金、校级优秀毕业生。}

%========================== 工作经历 ==========================
\section{工作经历}

\entryline{某科技公司\quad 高级软件工程师}{2025.07 -- 至今}
\item{负责排版引擎核心模块的开发与性能优化，主导重写词法分析层，
      吞吐量提升 2.1 倍。}
\item{搭建差分测试框架，覆盖 200+ 对照用例，线上缺陷率下降 60\%。}

\entryline{某研究院\quad 实习生}{2024.03 -- 2024.09}
\item{参与增量计算框架设计，实现脏区域追踪算法，重排延迟降低 45\%。}
\item{独立完成 GPU 渲染后端原型，支持无头纹理回读与亚像素抗锯齿。}

%========================== 项目经历 ==========================
\section{项目经历}

\entryline{NTex —— 现代 TeX 排版引擎}{2023.01 -- 至今}
\item{以 Rust 重写 TeX 宏求值器：Token 采用 8 字节 tagged union，
      宏体为连续不可变 TokenArray，并实现字节码双轨执行与等价性测试。}
\item{实现 DVI 逐字节一致的排版输出与 PDF 后端（Type1 字体嵌入）。}
\item{技术栈：Rust、wgpu/vello GPU 渲染、WASM 实时预览。}

\entryline{开源贡献}{2022.01 -- 至今}
\item{向某编译器项目提交 30+ PR，涉及寄存器分配与错误恢复。}

%========================== 专业技能 ==========================
\section{专业技能}

\item{{\bf 语言：}Rust（精通）、C/C++（熟练）、Python（熟练）。}
\item{{\bf 系统：}编译原理、增量计算、并行布局、GPU 渲染管线。}
\item{{\bf 工具：}Git、LLVM、perf、CI/CD 流水线搭建。}

%========================== 荣誉奖项 ==========================
\section{荣誉奖项}

\item{2024\quad 全国大学生编译器挑战赛 一等奖}
\item{2023\quad ACM-ICPC 亚洲区域赛 银牌}

\bye
